\documentclass[12pt]{article}
\usepackage{amsmath, amssymb}
\usepackage{graphicx}
\usepackage{geometry}
\geometry{margin=1in}

\title{AOE/CS/ME 6444 -- Verification and Validation in Scientific Computing \\
Homework \#2 -- Semester Project Definition}
\author{Cheng-Shun Chuang}
\date{Spring 2026}

\begin{document}

\maketitle

\section*{1. Application Area}

\textbf{Application:} Structural deflection of a residential wood header beam under uniform load.

The real system of interest is a wood header placed above a door opening in a residential building. The beam supports roof and floor loads and must satisfy serviceability limits on deflection.

For the purposes of this VV\&UQ project, the physical system will be simplified to a one-dimensional, steady Euler--Bernoulli beam problem.

\section*{2. Team Selection}

This project will be completed by Jason and Cheng-Shun.

\section*{3. System Description}

\subsection*{3.1 Real System}

The system consists of:

\begin{itemize}
\item A simply supported wood header beam
\item Span length $L$
\item Uniform distributed load $q$
\item Constant elastic modulus $E$
\item Second moment of area $I$
\end{itemize}

\subsection*{3.2 Surroundings and Environment}

The beam is supported by wall studs at each end. In reality, it is exposed to varying environmental loads and moisture effects. However, in this simplified model:

\begin{itemize}
\item Load is constant
\item Material properties are constant
\item No temperature or moisture coupling
\end{itemize}

\section*{4. Mathematical Model}

\subsection*{4.1 Governing Equation}

The steady Euler--Bernoulli beam equation is:

\begin{equation}
E I \frac{d^4 w(x)}{dx^4} = q
\end{equation}

where:

\begin{itemize}
\item $w(x)$ = vertical deflection
\item $x$ = spatial coordinate along the beam
\item $E$ = Young's modulus
\item $I$ = second moment of area
\item $q$ = uniform distributed load
\end{itemize}

This is a fourth-order ordinary differential equation (ODE) with one independent variable $x$.

\subsection*{4.2 Boundary Conditions}

For a simply supported beam:

\begin{align}
w(0) &= 0 \\
w(L) &= 0 \\
\frac{d^2 w}{dx^2}(0) &= 0 \\
\frac{d^2 w}{dx^2}(L) &= 0
\end{align}

These conditions enforce zero deflection and zero bending moment at the supports.

\section*{5. Analytical Solution}

For a simply supported beam under uniform load, the exact solution is:

\begin{equation}
w(x) = \frac{q}{24 E I}
\left(x^4 - 2Lx^3 + L^3 x\right)
\end{equation}

The maximum deflection occurs at midspan $x = L/2$:

\begin{equation}
\delta_{\max} = \frac{5 q L^4}{384 E I}
\end{equation}

This closed-form solution enables direct code verification.

\section*{6. Numerical Method}

The governing equation will be discretized using:

\begin{itemize}
\item Finite Difference Method (central difference approximation for $d^4w/dx^4$)
\item Uniform grid spacing $\Delta x$
\end{itemize}

Mesh refinement studies will be performed to determine:

\begin{itemize}
\item Observed order of accuracy
\item Grid convergence
\item Discretization error
\end{itemize}

\section*{7. Verification Plan}

\subsection*{7.1 Code Verification}

\begin{itemize}
\item Compare numerical solution with analytical solution.
\item Compute $L_2$ error norm.
\item Perform systematic mesh refinement.
\item Compute observed order of accuracy.
\item Implement Method of Manufactured Solutions (MMS).
\end{itemize}

\section*{8. Validation Plan}

Validation will compare numerical predictions against:

\begin{itemize}
\item Classical beam theory solutions
\item Residential building code serviceability limits (e.g., $L/240$ deflection limit)
\end{itemize}

\section*{9. Uncertainty Quantification}

Uncertainty will be introduced in:

\begin{itemize}
\item Load $q \sim \mathcal{N}(\mu_q, \sigma_q)$
\item Modulus $E \sim \mathcal{N}(\mu_E, \sigma_E)$
\end{itemize}

Monte Carlo simulation (100--1000 samples) will be used to compute:

\begin{equation}
P(\delta_{\max} > L/240)
\end{equation}

\section*{10. Summary}

This project focuses on verification, validation, and uncertainty quantification of a one-dimensional structural beam model representing residential header deflection. The governing equation is a fourth-order ODE with a known analytical solution, making it well-suited for rigorous VV\&UQ analysis.

\end{document}